\section{Technological Structure}
\label{sec:model}

\subsection{Notation and Basic Definitions}

Consider an economy with $n$ production sectors indexed by
$i,j = 1,\ldots,n$.
Following \citet{leontief1986input}, we adopt the following notation.

\begin{definition}[Economic Variables]
At each time period $t$, the economy is characterized by:
\begin{itemize}
    \item $\vec{X}_t \in \mathbb{R}^n_+$: gross output vector, where
          $X_{i,t}$ is total production in sector $i$;
    \item $\vec{C}_t \in \mathbb{R}^n_+$: realized household
          consumption vector;
    \item $\hat{\vec{C}}_t \in \mathbb{R}^n_+$: revealed (estimated)
          consumption demand;
    \item $\vec{I}_t \in \mathbb{R}^n_+$: gross investment vector;
    \item $\vec{K}_t \in \mathbb{R}^n_+$: capital stock vector.
    \item $\vec{l} \in \mathbb{R}^n_+$: labour coefficient vector,
          where $l_j$ is labour required per unit of output in sector
          $j$;
    \item $L$: total available labour in the economy;
\end{itemize}
\end{definition}

\begin{definition}[Technology Matrices]
The technological structure is described by:
\begin{itemize}
    \item $A \in \mathbb{R}^{n \times n}_+$: intermediate input
          coefficient matrix, where $A_{ij}$ is the amount of good $i$
          required per unit of output in sector $j$;
    \item $B \in \mathbb{R}^{n \times n}_+$: capital coefficient
          matrix, where $B_{ij}$ is the capital of type $i$ needed per
          unit of output in sector $j$;
    
    \end{itemize}
\end{definition}

\subsection{Material Balance and Capital Accumulation}

Following the Leontief system \citep{leontief1986input}, the material
balance equation states that net output equals final demand:
%
\begin{equation}\label{eq:material_balance}
    (I - A)\vec{X}_t = \vec{C}_t + \vec{I}_t + \vec{G}_{t}
\end{equation}
%
where $I$ is the $n \times n$ identity matrix.
The matrix $(I - A)$ must satisfy $\rho(A) < 1$ for the system to be
productive \citep{nikaido1970introduction}.

Capital accumulation follows \citet{von1945model}:
%
\begin{equation}\label{eq:capital_accumulation}
    \vec{K}_{t+1} = (1 - \delta)\vec{K}_t + \vec{I}_t,
\end{equation}
%
where $\delta \in (0,1)$ is the depreciation rate.

The capital feasibility condition ties the capital stock to gross
output:
%
\begin{equation}\label{eq:capital_constraint}
    \vec{K}_{t} \geq B\vec{X}_{t}.
\end{equation}
%
Here $B_{ij}$ gives the stock of capital good $i$ required to produce
one unit of output in sector $j$.
This formulation imposes constant returns to scale on capital usage,
which is standard in input-output theory.

\subsection{Maintenance Investment}

At depreciation rate $\delta$, the investment required merely to
maintain the capital stock is:
%
\begin{equation}\label{eq:maintenance}
    \vec{I}^{\mathrm{maint}}_t = \delta \vec{K}_t = \delta B \vec{X}_t,
\end{equation}
%
where the second equality uses the binding capital constraint
$\vec{K}_t = B\vec{X}_t$.
Substituting into~\eqref{eq:material_balance}:
%
\begin{equation}\label{eq:modified_balance}
    (I - A - \delta B)\vec{X}_t = \vec{C}_t +\vec{G}_{t}+ \Delta \vec{K}_t,
\end{equation}
%
where $\Delta \vec{K}_t = \vec{K}_{t+1} - \vec{K}_t$ is net capital
formation.

\subsection{Augmented Technical Matrix}

\begin{definition}[Augmented Technical Coefficients]
\label{def:augmented}
Define the augmented technical coefficient matrix:
\begin{equation}\label{eq:A_bar}
    \bar{A} = A + \delta B.
\end{equation}
\end{definition}
%
Then equation~\eqref{eq:modified_balance} becomes:
%
\begin{equation}\label{eq:augmented_balance}
    (I - \bar{A})\vec{X}_t = \vec{C}_t + \vec{G}_{t} + \Delta \vec{K}_t.
\end{equation}

\begin{assumption}[Productivity of Augmented System]
\label{ass:productivity}
The augmented system is productive: $\rho(\bar{A}) < 1$, ensuring that
$(I - \bar{A})^{-1}$ exists and has non-negative entries
\citep{nikaido1970introduction}.
\end{assumption}

\subsection{Discrete-Time Capital Dynamics and Dynamic Multiplier}

Beginning from the augmented material balance~\eqref{eq:augmented_balance},
which holds at every period $t$:
%
\begin{equation}
    (I - \bar{A})\vec{X}_{t} = \vec{C}_{t} + \vec{G}_{t} + \Delta\vec{K}_{t},
\end{equation}
%
where $\Delta\vec{K}_{t} = \vec{K}_{t+1} - \vec{K}_{t}$ is net capital
formation. Since the augmented system is productive
(Assumption~\ref{ass:productivity}), the Leontief inverse
$(I - \bar{A})^{-1}$ exists and is non-negative, so gross output is
determined by final demand:
%
\begin{equation}
    \vec{X}_{t} = (I - \bar{A})^{-1} \vec{F}_{t},
    \qquad
    \vec{F}_{t} \equiv \vec{C}_{t} + \vec{G}_{t} + \Delta\vec{K}_{t}.
\end{equation}
%
Since the capital constraint $\vec{K}_{t} = B\vec{X}_{t}$ holds with
equality at every period, it also holds one period forward, so that:
%
\begin{equation}
    \Delta\vec{K}_{t} = B\,\Delta\vec{X}_{t}.
\end{equation}
%
Because $(I - \bar{A})^{-1}$ is a constant matrix, first-differencing
the material balance gives $\Delta\vec{X}_{t} = (I-\bar{A})^{-1}
\Delta\vec{F}_{t}$, and therefore:
%
\begin{equation}\label{eq:delta_K_first}
    \Delta\vec{K}_{t} = M\,\Delta\vec{F}_{t},
\end{equation}
%
where the \emph{dynamic multiplier matrix} is:
%
\begin{equation}\label{eq:multiplier}
    M = B(I - \bar{A})^{-1}.
\end{equation}
%
Expanding $\Delta\vec{F}_{t} = \Delta(\vec{C}_{t} + \vec{G}_{t}) +
\Delta^{2}\vec{K}_{t}$ and substituting into~\eqref{eq:delta_K_first}:
%
\begin{equation}
    \Delta\vec{K}_{t}
    = M\,\Delta(\vec{C}_{t} + \vec{G}_{t})
    + M\,\Delta^{2}\vec{K}_{t}.
\end{equation}
%
Since relation~\eqref{eq:delta_K_first} holds at every period, it
holds for all higher-order differences as well:
$\Delta^{k}\vec{K}_{t} = M\,\Delta^{k}\vec{F}_{t}$ for each
$k \geq 1$. Applying this recursively to the remainder term at each
order yields, after $K$ iterations:
%
\begin{equation}
    \Delta\vec{K}_{t}
    = \sum_{k=1}^{K} M^{k}\,\Delta^{k}(\vec{C}_{t} + \vec{G}_{t})
    + M^{K}\,\Delta^{K+1}\vec{K}_{t}.
\end{equation}
%
The remainder $M^{K}\Delta^{K+1}\vec{K}_{t} \to \vec{0}$ as
$K\to\infty$ provided $\rho(M) < 1$, which is implied by the
productivity of the augmented system. Taking the limit gives the
exact discrete Neumann series:
%
\begin{equation}\label{eq:finite_difference}
\boxed{
    \Delta\vec{K}_{t}
    = \sum_{k=1}^{\infty} M^{k}
    \bigl(\Delta^{k}\vec{C}_{t} + \Delta^{k}\vec{G}_{t}\bigr),}
\end{equation}
%
which converges provided $g \cdot \rho(M) < 1$. Here, $g$ is the homogenous rate at which $\vec{F}_t$ grows.

\begin{remark}
    Equation~\eqref{eq:finite_difference} is the standard result of
    the dynamic Leontief model with the depreciation-augmented
    technical matrix $\bar{A} = A + \delta B$ in place of $A$.  The
    derivation requires only that the augmented material balance and
    the capital constraint hold with equality at every discrete period;
    no continuous-time approximation is introduced.
\end{remark}

\begin{remark}
    Because household preferences evolve through the adaptive updating
    rule described in Section~\ref{sec:dynamics}, the future utility
    function is unknown at time $t$. A conventional Euler equation for
    optimal investment therefore cannot be derived. Instead, we
    introduce a feedback investment rule motivated
    by~\eqref{eq:finite_difference}; see Section~\ref{sec:growth}.
\end{remark}